\documentclass{article}
\usepackage{amsmath,amssymb,mathtools}
\newtheorem{lemma}{Lemma}
% The expected output of `leanify check --claims` is check.txt: composition,
% inclusion, suprema and infima, finite sums and products, equivalent statements,
% `$$` displays, abbreviations, standard functions and constants, set operations,
% connectives and quantifiers; sentences of several parts, the positive part, and a
% symbol of a result beside a notation of the article with the same letter.
\begin{document}
Define $f\colon(0,\infty)\to\mathbb{R}$ by $f(x)\coloneqq x+\frac1x$.
Define $g\colon\mathbb{R}\to\mathbb{R}$ by $g(x)\coloneqq x^2+1$.

\begin{lemma}\label{lem:comp}
For every $x\in\mathbb{R}$, $(f\circ g)(x)\ge2$.
\end{lemma}

\begin{lemma}\label{lem:sub}
$(0,1)\subseteq[0,1]$.
\end{lemma}

\begin{lemma}\label{lem:sup}
$[0,1]\supseteq(0,1]$.
\end{lemma}

\begin{lemma}\label{lem:bounds}
$\sup_{x\in(0,1)} x^2=1$.
\end{lemma}

\begin{lemma}\label{lem:inf}
$\inf\{f(x) : x\in(0,\infty)\}=2$.
\end{lemma}

\begin{lemma}\label{lem:sums}
$\sum_{j=1}^{4} j=10$.
\end{lemma}

\begin{lemma}\label{lem:products}
For every $x\in\mathbb{R}$, $$\prod_{k=1}^{2}(x+k)=x^2+3x+2.$$
\end{lemma}

\begin{lemma}\label{lem:tfae}
Let $x\in\mathbb{R}$.
Then the following are equivalent:
\begin{enumerate}
\item $x^2=1$;
\item $\lvert x\rvert=1$;
\item[(c)] $(x-1)(x+1)=0$.
\end{enumerate}
\end{lemma}

\begin{lemma}\label{lem:functions}
For every $x\in(0,\infty)$, $\ln x\le x-1$.
Moreover, for every $y\in\mathbb{R}$, $\cos^2 y+\sin^2 y=1$.
\end{lemma}

\begin{lemma}\label{lem:exp}
For every $x\in\mathbb{R}$, $e^{x}\ge1+x$ and $\exp(x)>0$.
\end{lemma}

\begin{lemma}\label{lem:pi}
$3<\pi$ and $\pi<\frac{22}{7}$.
\end{lemma}

\begin{lemma}\label{lem:setops}
$[0,1]\cap[1,2]=\{1\}$ and $[0,2]\setminus(0,2)=\{0,2\}$.
\end{lemma}

\begin{lemma}\label{lem:empty}
For every $x\in\mathbb{R}_{>0}$, $x\notin\emptyset$ and $x\in\{y\in\mathbb{R}\mid y>0\}$.
\end{lemma}

\begin{lemma}\label{lem:logic}
$\forall x\in\mathbb{R}: \exists y\in\mathbb{R}: y>x$.
Moreover, for every $x\in\mathbb{R}$, $x>1\Rightarrow x^2>1$ and $\neg(x<x)$.
\end{lemma}

\begin{lemma}\label{lem:misc}
For every $x\in\mathbb{R}$, $\max(x,-x)=|x|$ and $\operatorname{sgn}(x)\,x=|x|$. Moreover, $\sqrt[3]{8}=2$, $2\times3=6$ and $\int_0^1 x\,dx=\frac12$.
\end{lemma}

\begin{lemma}\label{lem:conditions}
For every $x\in\mathbb{R}$ with $x\neq0$, $x^{-1}x=1$.
Moreover, for every $y\in\mathbb{R}\setminus\{0\}$, $y^{-2}>0$.
\end{lemma}

\begin{lemma}\label{lem:integers}
For every $x\in\mathbb{R}$, $\lfloor x\rfloor\le x<\lfloor x\rfloor+1$ and $x\le\lceil x\rceil$.
Moreover, for every $n\in\mathbb{N}_0$ with $n\ge1$, $\frac1n\le1$.
Finally, $\mathbb{N}\subseteq\mathbb{Z}$.
\end{lemma}

\begin{lemma}\label{lem:finite-sums}
For every $n\in\mathbb{N}$, $\sum_{j=1}^{n} j=\frac{n(n+1)}{2}$.
Moreover, for every $n\in\mathbb{N}$ with $n\ge1$, $\prod_{k=1}^{n} k=n!$.
\end{lemma}

\begin{lemma}\label{lem:series}
$\sum_{k=0}^{\infty}\frac{1}{2^k}=2$ and $\sum_{k=1}^{\infty}\frac{(-1)^{k+1}}{k}=\ln 2$.
\end{lemma}

\begin{lemma}\label{lem:limits}
$\lim_{n\to\infty}\frac1n=0$, $\lim_{x\to0}\frac{\sin x}{x}=1$ and $\lim_{x\to\infty} e^{-x}=0$.
\end{lemma}

\begin{lemma}\label{lem:binomial}
For all $n,k\in\mathbb{N}$ with $k\le n$, $\binom{n}{k}=\frac{n!}{k!\,(n-k)!}$.
\end{lemma}

The value of $f$ at $1$, i.e.\ $f(1)$, is $2$, cf.\ Lemma~\ref{lem:inf}.

The defining sign is read in the ways articles write it.
Set $c_1:=2$, $c_2\triangleq c_1+1$ and $c_3\vcentcolon=c_2^2$.
For $x\in\mathbb{R}$, set $q(x)\mathrel{:=}x^2+c_1$.

\begin{lemma}\label{lem:defining-signs}
$c_3=9$ and $q(1)=3$.
\end{lemma}

\section{The words and the integrals of the core}

The words about numbers that every article uses need no field: \emph{real number}, \emph{integer}, \emph{positive}, \emph{rational number}.
An integral $\int_a^b$ is the oriented integral of a function of a real variable; it is meaningful when the function is integrable between $a$ and $b$.

\begin{lemma}\label{lem:core-words}
For all real numbers $x$ and $y$, $x^2+y^2\ge2xy$.
Moreover, for every positive integer $n$ and every nonnegative real number $a$, $a^n\ge0$.
There is no rational number $q$ such that $q^2=2$.
\end{lemma}

\begin{lemma}\label{lem:integrals}
We have $\int_1^2\frac1x\,dx=\log 2$.
Moreover, for every $a>0$ and every $b>a$, $\int_a^b x\,dx=\frac{b^2-a^2}{2}$.
Finally, $\int_0^1\int_0^1 st\,ds\,dt=\frac14$.
\end{lemma}

An integrand has to be defined at almost every point between the ends, not at every point.

\begin{lemma}\label{lem:integral-ae}
We have $\int_0^1\frac{1}{\sqrt x}\,dx=2$.
\end{lemma}

An integral may be over a ray or over the line.

\begin{lemma}\label{lem:integral-improper}
We have $\int_{-\infty}^\infty e^{-t^2}\,dt=\sqrt\pi$ and $\int_1^\infty\frac1{t^2}\,dt=1$.
Moreover, for every $x\in\mathbb{R}$, $\int_{-\infty}^x e^{t}\,dt=e^x$.
\end{lemma}

A variable whose domain uses the variables before it brings them along: a quantity defined with it depends on them too.

\begin{lemma}\label{lem:dependent-domain}
Let $n\ge1$ and $x\in[0,n]$.
Set $r\coloneqq\frac{1}{x+1}$.
Then $r\le1$.
Moreover, for every $t\in(0,1)$, $r\le1+t$.
\end{lemma}

The same holds for a family: the domain of its argument may use its index.
For $c\in(0,2)$ and $t\in[0,c]$, let $\kappa_c(t)\coloneqq c-t$.

\begin{lemma}\label{lem:dependent-family}
We have $\kappa_1(\frac12)=\frac12$.
\end{lemma}

The indicator of a condition is $1$ where the condition holds and $0$ elsewhere, and the indicator of a set is that of membership.
For $v\in[0,1]$ and $x\in\mathbb{R}$, set $e_v(x)\coloneqq\mathbf 1\{x\le v\}-v$.
For $x\in\mathbb{R}$, set $w(x)\coloneqq\mathbf{1}_{[0,1]}(x)+\mathbb{1}_{\{1<x,\ x\le2\}}$.

\begin{lemma}\label{lem:indicator}
For every $v\in[0,1]$, $\int_0^1e_v(x)\,dx=0$.
Moreover, for every $x\in\mathbb{R}$, $0\le w(x)\le1$.
Finally, $\int_0^2w(x)\,dx=2$.
\end{lemma}

A family written with a subscript does not take its letter: beside such a family, the
letter is still Euler's number in a power.

\begin{lemma}\label{lem:euler}
We have $e^0=1$.
Moreover, for every $v\in[0,1]$, $e_v(0)\le e^v$.
\end{lemma}

We write $h_a$ for $a^2+1$.
For $x\in\mathbb{R}$, set $\phi(x)\coloneqq x^2+1$, and let $\psi(x)\coloneqq2\phi(x)$.
We write $s(t)\coloneqq t^3-t$, and $a\star b$ for $ab+1$.
For $x\in\mathbb{R}$, set $\theta(x)\coloneqq(x-1)^+-(x-1)^-$.
We consider two sums of a positive integer $n$,
\[
S_1(n)\coloneqq\sum_{k=1}^{n}k,\qquad S_2(n)\coloneqq\sum_{k=1}^{n}k^2.
\]

\begin{lemma}\label{lem:parts}
For every $x\in\mathbb{R}$, $\psi(x)\ge2$, $x\star0=1$ and $\theta(x)=x-1$.
Moreover, $s(1)=0$ and $S_1(2)=3$.
\end{lemma}

\begin{lemma}\label{lem:local}
Let $y>0$. For $x\in\mathbb{R}$, let $h_y(x)\coloneqq x+y$. Then $h_y(0)=y$.
\end{lemma}

\begin{lemma}\label{lem:family}
For $t\in[0,1]$ and $x,y\in\mathbb{R}$, let $u_t(x)\coloneqq tx$ and $v_t(x,y)\coloneqq u_t(x)+y$.
Then $v_t(0,0)=0$.
\end{lemma}

\begin{lemma}\label{lem:notation}
$h_2=5$.
\end{lemma}
\begin{lemma}\label{nested-bars}
For all real numbers $x$ and $y$, $||x|-|y||\le|x-y|$ and $\big||x|-|y|\big|\le|x|+|y|$. For all real numbers $a$ and $b$ with $a<b$, $a<\frac{a+b}{2}<b$. For all real numbers $a$, $b$ and $c$ with $a<b<c$, $a<c$.
\end{lemma}

\end{document}
